\chapter{Segmentation and the threshold sweep}
\label{ch:segment}

\texttt{AMRSegment} finds the connected components of a threshold set, and
\texttt{amrSweep} walks the threshold and reports how they behave. Together they
are the substrate for \ref{req:compare}: the smallest apparatus that produces a
scientific result about candidate cell definitions.

\section{Components of a threshold set}

Components of $\sublev{v}$---or of $\{f < v\}$---are found by union--find over
the complex's own adjacency, which already carries the periodic wrap and the
coarse--fine stitch. No boundary plane enters, so \ref{req:periodic} holds with
nothing further required.

\paragraph{Labels are the minimum identifier in the component.}
This is a property of the component itself and not of who found it, so the same
run at any rank count produces identical labels. Segmentations are therefore
directly comparable across rank counts, across time snapshots, and across
different candidate fields---which is precisely what \ref{req:compare}
and~\ref{req:track} require.

\paragraph{The cross-rank merge exploits the same fact.}
Each rank publishes $(\gid, \text{local root})$ pairs; two ranks naming the same
$\gid$ with different roots is exactly the information needed to join their
components, and every rank replays the same union--find over the gathered pairs,
so no broadcast is needed.

Only nodes another rank might also see are published: ghosts, owned nodes within
one cell of a box boundary, and---the case easy to miss---owned nodes adjacent to
a covered cell, since those may be stamped as coarse parents into a finer box's
ghost region and so be seen by whichever rank owns that fine box. Omitting the
third case would fail to join components across a coarse--fine interface,
silently and only under MPI. Publishing every node instead, as the first version
did, is $O(N)$ all-to-all: it passes at four ranks on a $64^3$ box and fails on
the snapshots the tool exists for.

\section{Volumes are integer-derived}
\label{sec:intvol}

Cell volumes take exactly one value per level. The tally therefore counts
integers per (component, level) and multiplies at the end:
\begin{equation}
  V_c \;=\; \sum_{\lev} n_{c,\lev}\,\dvol_{\lev},
  \qquad \dvol_{\lev} = \prod_{k=1}^{\dsp} \Delta x_{\lev,k}.
  \label{eq:volume}
\end{equation}
Because uncovered cells partition the domain, \cref{eq:volume} is exact, and
because it is a function of integers it is bit-identical at any decomposition.

\begin{designrule}\label{rule:intout}
Where a scientific output can be derived from integers, derive it from integers.
A quantity that moves in its last bits when the rank count changes is
comparable only to a tolerance, and here that was avoidable rather than
inherent.
\end{designrule}

\section{The sweep}

\texttt{amrSweep} walks the threshold over a linear or logarithmic range and
reports, at each value: the number of connected regions, the total volume, the
largest and median component volume, and the isosurface area. Run on several
candidate surrogates, the tables are directly comparable.

What to look for is a plateau in $N(v)$, or a gap in the size distribution.
Either is a threshold-independent scale and therefore a defensible cell
definition. A featureless continuum is also an answer, and arguably the more
interesting one, because it says that no threshold is defensible and shifts the
question.

The sweep is a coarse sampling of a merge tree, which makes it a free and sharp
regression gate:

\begin{designrule}\label{rule:sweepgate}
When the merge tree is built it must reproduce \texttt{amrSweep}'s $N(v)$
exactly at every sampled threshold.
\end{designrule}

\section{Verification against closed forms}

\texttt{makeTestPlt} writes a cellular field
\begin{equation}
  f(\mathbf{x}) \;=\; \sum_{k=1}^{\dsp} \cos\!\left(\frac{2\pi n\,x_k}{L_k}\right),
  \label{eq:cellfield}
\end{equation}
which on a periodic domain has exactly $n^{\dsp}$ maxima, and saddles between
neighbours at a value that can be written down: moving between adjacent maxima
one term runs from $+1$ to $-1$ while the others stay at $+1$, so the saddle sits
at $\dsp - 2$.

\begin{table}[htbp]
\centering
\caption{Component count against the analytic saddle of \cref{eq:cellfield},
  $n=3$. The maxima lie \emph{on} the domain faces, so recovering exactly
  $n^{\dsp}$ of them is also a test that the periodic wrap is present in the
  adjacency rather than assumed.}
\label{tab:cellcounts}
\begin{tabular}{llll}
\toprule
case & above the saddle & below & analytic saddle \\
\midrule
3D, $n=3$ & 27 components & 1 & exactly $1$ \\
2D, $n=3$ &  9 components & 1 & exactly $0$ \\
\bottomrule
\end{tabular}
\end{table}

A threshold sitting exactly \emph{on} a saddle value gives a genuinely
ill-defined count---the 3D case reads $20$ at $v = 1.000$. That is the
degeneracy persistence exists to quantify, not a defect.

The sphere case checks volume and area against closed forms to $0.06\%$ and
$0.04\%$. Counts and volumes are unchanged across \texttt{max\_grid\_size}
$8/16/32/64$ and across a three-level hierarchy.
